\documentclass[11pt,a4paper]{article}
\usepackage[T1]{fontenc}
\usepackage[utf8]{inputenc}
\usepackage[english]{babel}
\usepackage{amsmath,amssymb}
\usepackage[margin=2cm]{geometry}
\setlength{\parindent}{0pt}
\setlength{\emergencystretch}{3em}
\allowdisplaybreaks

\title{About \texttt{worked\_derivatives.tex}}
\author{}
\date{}
\begin{document}
\maketitle
This note explains what the test data in \texttt{worked\_derivatives.tex} is used for. That file is input for the engine in \texttt{main.c}; it is read by the test \texttt{test\_worked\_derivatives.c} in this folder.

\begin{enumerate}
\item \textbf{Purpose.} The file holds 6 requests $\frac{d}{dx}(\ldots)$. The answers are checked elsewhere. Here the engine must also record the way: one line for each piece of the expression that contains $x$, from the innermost outwards, with the derivative of the piece and the rule that gives it.

\item \textbf{What it exercises.} Rule 27n of \texttt{main.c} and the sentences of chapter 41 of \texttt{output.c}. Every line is checked on its own before it is recorded: the piece and its derivative are read back from the text that will be printed and sent to \texttt{checker.c}, which differentiates the piece with rules of its own.

\item \textbf{Variety.} The power rule, a constant factor, sums, the product rule, the quotient rule, the exponential, a sine, a tangent, a square root, and the chain rule with a first-degree, a second-degree and a power inside.

\item \textbf{Expected result.} The line \texttt{PASS: 124 of 124 checks on worked derivatives passed}. The expected lines are written in the test as formulas worked out by hand, and compared with the recorded ones by value, not by spelling.

\item \textbf{How to run it.} From the top folder of the repository:
\begin{verbatim}
cc -std=c11 -O2 -Wall -Wextra -o /tmp/test_worked_derivatives \
    tests/derivatives/test_worked_derivatives.c \
    && /tmp/test_worked_derivatives
\end{verbatim}
\end{enumerate}
\end{document}
